\section{Ablations and Instrument Findings}
\label{sec:ablations}

\subsection{Static vs.\ active restricted access}
\label{sec:mechanism}
All unsupported commitments arise on the \emph{active} strict condition; the
static restricted condition is $10/10$ correct for both models. A partially
blind coding of GPT's seven strict accepts (codebook frozen before coding;
packets masked D01--D07) assigns the primary code \textbf{C3}---treating
\emph{absence of failure evidence as evidence of absence}---to all seven,
co-firing C1 (symptom resolution), C2 (coverage completeness), and C5
(provenance/freshness); \textbf{C4} (explicitly acknowledging the withholding)
never fires. This is a \emph{single AI rule-based coder with no independent
human coder}, so it is a qualitative hypothesis only. Moreover activeness,
freshness, and trust changed together across protocol versions, so the
active/static comparison is confounded; an activeness-only controlled arm
(freshness and trust held fixed) is required before calling this a demonstrated
mechanism.

\subsection{Instrument residual: baseline P2P projection gap}
The tier-1 non-interference scan was $61/61$ clean, but a tier-2 minimisation
found the \emph{pre-patch} \texttt{pass\_to\_pass\_baseline} per-test map
surviving the tier-D projection (all \textsc{Pass}, zero \textsc{Fail}). It
cannot reconstruct gold (the post-patch P2P map is removed), was reachable only
by explicit fetch, and was seen by none of the seven error cells: the single GPT
cell that saw it (20590) abstained, and DeepSeek saw it in eight cells and
abstained throughout. It therefore has no causal link to the reported errors,
but it is a disclosed non-interference imperfection and is mandated for removal
before the next live run. The corrected v2 projection used for the prepared
80-cell validation bed removes this field on every tool path
(Section~\ref{sec:validation}).

\subsection{Independent-issue pilot status}
The eight G3-C.1 issues rest on 40 gold certificates, evidence projections
with $144/144$ clean restricted surfaces (224 surfaces total), $456/456$ HTTP
cross-checks, and frozen manifests. This measurement bed was subsequently
exercised by the live 80-cell run (Section~\ref{sec:e80}). Mock cells used
during engineering contribute \emph{zero} model-performance datapoints, are
not used to fill any results cell, and are excluded from research analysis;
every reported datapoint comes from a real, raw-first-persisted model
response.

\subsection{External validation: engineering readiness and live execution}
\label{sec:validation}
An 80-cell external validation on the eight new issues (2 models $\times$
5 conditions $\times$ 8 issues) was engineered, cleared every offline gate at
\$0, received separate written authorization, and then ran to completion
(Section~\ref{sec:e80}). Pre-paid readiness checks reported:
\begin{itemize}[leftmargin=1.6em,itemsep=1pt,topsep=2pt]
  \item \textbf{non-interference:} 224 surfaces probed; 144 restricted
  tier-D; \textbf{0} tier-1, \textbf{0} tier-2, \textbf{0} baseline, and
  \textbf{0} hidden-gold exposures; $40/40$ dedicated mock gold alignments
  ($64/64$ non-Escalate cells through the generic script);
  \item \textbf{protocol:} strict 403 on all 8 restricted fetches; escalation
  400-without-reason then 200-with-reason on all 8; v2 projection tag present;
  \item \textbf{model id:} exact served-id match, hard stop on drift, no false
  trip when id absent;
  \item \textbf{budget/resume/raw-first:} guards at \$2.00 total / \$1.00 per
  model (allow + per-cell cap + per-model block + total block, both models);
  resume restores accumulated spend and refuses NaN/Inf/negative ledgers;
  raw-first persistence sha-verified on all 80 raw artifacts before scoring.
\end{itemize}
The pre-run engineering verdict was \texttt{ENGINEERING\_READY}, with three
\emph{non-engineering} items tracked in
\path{paper_sprint/experiment80/READINESS.md}: (i)~\textbf{authorization}
for live model calls, since granted in writing for 80 cells at the
\$2.00/\$1.00 caps; (ii)~\textbf{gold authority}---\texttt{annotator2 =
NOT\_OBTAINED}, so the run is reported as a \emph{Provisional-Gold Exploratory
Evaluation}, not human-arbitrated benchmark validation (PROVISIONAL by
design); and (iii)~\textbf{runner binding}---a dedicated live runner
(\texttt{g3c1live-runner-v1}), instantiated for the eight G3-C.1 issues at
\$2.00/\$1.00 and fingerprinted pre-run, rather than the frozen G3-B.3
entrypoint (which defaults to \$3.00/\$1.50 and the G3-A store). The bound
runner completed $80/80$ cells, 162 calls, spend \$0.258198, with zero
retries, terminal errors, response-format downgrades, and resume-skips.

\subsection{From instrument findings to a controlled intervention}
The mechanism codes above (notably C2/C5: completeness and provenance acting
as acceptance warrants despite withheld decisive results) are associational
and rely on partially blind qualitative coding. G4
(Section~\ref{sec:g4}) tests that mechanism causally: linkage is fixed, the
evidence transcript is yoked, and only decisive-outcome visibility and
freshness/coverage attestation vary. The G4 H1 result supplies the
controlled counterpart to the C5-style explanation suggested here.
